\documentclass[12pt]{article}
\usepackage[utf8]{inputenc}
\usepackage[T1]{fontenc}
\usepackage{hyperref}
\title{Yamap Language Resources for K-12 Curriculum Development}
\author{Compiled by Claude Sonnet 5.5 (claude-sonnet-5-5) for Andr\'as Kornai}
\date{}
\begin{document}
\maketitle

\section{Introduction}
Yamap (ISO 639-3: ymp; Glottocode yama1260) is an Oceanic Austronesian language of Morobe Province, Papua New Guinea. It belongs to the Hote subgroup of the Huon Gulf languages. Wikipedia cites about 1,600 speakers (2000 census). Global Recordings gives a population of 670 and the alternate name ``Ombe''. Joshua Project reports roughly 3,200 people in the group. Yamap is unrelated to the Papuan ``Yam'' languages of the Morehead River area. The searches found no school materials and no dictionary.

\section{Primary Educational Resources}
No primers, textbooks or school programmes appeared in the searches.

\section{Linguistic and Reference Materials}
\begin{itemize}
\item Wikipedia, ``Yamap language''. According to its external links, a Paradisec open-access collection of materials from Meinrad Scheller exists. I did not verify what that collection contains. \url{https://en.wikipedia.org/wiki/Yamap_language}
\end{itemize}

\section{Religious and Community Resources}
\begin{itemize}
\item Joshua Project profile of the Yamap people (demographic and mission-oriented summary): \url{https://joshuaproject.net/people_groups/15970}
\end{itemize}

\section{Digital and Audio Resources}
\begin{itemize}
\item Global Recordings Network language page for ymp. It reports no audio recordings in Yamap. \url{https://globalrecordings.net/en/language/ymp}
\end{itemize}

\section{Curriculum Development Context}
Public resources are very scarce. The most promising lead is the Paradisec collection named on Wikipedia, which should be checked directly for word lists, recordings and texts. Since Yamap is spoken by all members of its community, a native-language curriculum would rest on community knowledge and new documentation.
\end{document}